\section{Additional Artifact Details}
\label{app:artifacts}

\subsection{Corpus accounting}

The plan corpus is the full cross product of ten topologies, five intent
classes, two prompt variants (invariants stated or withheld) and five
repetitions: \(10\times5\times2\times5=500\) cells, one agent call each.
Generation uses \texttt{qwen2.5:7b-instruct-q8\_0} at temperature $0.7$ with
a 16{,}384-token context. The sampling seed of a cell is derived from the
cell's own identity rather than from execution order, so the corpus is
reproducible and independent of scheduling. Each record stores the prompt
hash and length, the raw response, the seed, prompt and generation token
counts, the parsed plan, per-class ground-truth violation labels under both
temporal scopes and all three budgets, the intent-fulfillment verdict, and
the a-priori feasibility label the generator recorded. Certified feasibility
is a separate offline verdict, stored beside the corpus and joined at read
time rather than written back into it: the corpus records what was asked and
what came back, and a corpus that changes under reanalysis is not evidence.

\subsection{Cost sweep accounting}

The sweep covers 26 topologies $\times$ 2 backends $\times$ 2 temporal
scopes $\times$ 3 budgets $=312$ cells. Each cell times six clean plans and
six violating plans, five repetitions each, under three measurement
conditions---clean full traversal, violating with early termination, and
violating with early termination disabled---for \(28{,}080\) timed
verifications in 3{,}045\,s. Reported values are the mean across plans of
each plan's per-repetition median.

Three conditions are separated because they answer different questions. A
plan that is rejected at its first violated predicate is cheap to check, and
timing only such plans would credit a profile for detecting early, which is
a detection property rather than a cost. Timing only clean plans has the
opposite bias for this encoder: its configuration-integrity fragment
instantiates constraints only where a conflict concretely exists, so depth
three costs exactly what depth two costs on a clean plan. Assumption
monotonicity is therefore asserted on the third condition, the deterministic
worst-case envelope, which is the only series measuring the same work at
every budget. It would be convenient to settle the deepest step by counting
emitted disjuncts instead of by timing---the integrity layer instantiates a
constraint only where a conflict concretely exists, and on a clean state
emits none---but the envelope is the maximum over clean and violating plans,
and on a violating state that layer does emit; depth three additionally runs
a structural pass over the whole plan that shallower depths skip.
Monotonicity is reported as an assumption held to the measurement, not one
the encoding closes.
Stability is judged per cell by the dispersion of repeated timings rather
than by machine load, which on a workstation is dominated by unrelated
processes; one cell of 312 exceeded the threshold and is named in the frozen
table.

\subsection{Closed-loop sweep accounting}

One hundred intents---five topologies, five intent classes, two prompt
variants, two repetitions---are each run under eight policies, giving 800
traces. Every policy draws its round-zero plan with the same seed, so all
eight begin from an identical plan; this was verified to hold in all 100
cells. Each trace stores every round's prompt, plan, verifier verdict,
counterexample, verification time and agent latency, and---on actuation---the
pre- and post-state hashes together with the violations found by rechecking
the state that was actually realized.

Paired tests use the intent as the unit. Reporting each policy's 100 traces
as independent observations would treat eight views of the same intent as
eight samples; the exact McNemar tests in Table~II of the paper instead
count the intents on which two policies disagree.

\subsection{Selection protocol}

The profile space searched is the cross product of two temporal scopes and
three budgets on the SMT backend: six profiles. The twin backend is not a
candidate, because the agent corpus yields no per-topology detection estimate
for it; the reported selection result is therefore over coverage depth and
temporal scope, not over the backend dimension. For each of ten held-out
topologies and each of the thirteen risk prices in the sweep, primitives are
pooled over the other nine topologies, the profile minimizing
\(J=R\,V+\E[T]\) is chosen, and that profile is scored with the held-out
topology's own \(q\), \(p\), \(\tau\) and \(G\): 130 selections in total.
The reference is a plug-in oracle---the profile minimizing the same closed
form under the held-out topology's own primitives---not a closed-loop
optimum. All four primitives are estimated on one population, certified-%
feasible intents whose plan fulfills them, so that they can be combined. The
corpus is rescored offline at every profile, so the experiment adds no agent
calls. Detection is recomputed by running the verifiers over the stored plans
rather than read from stored labels, so numerator and denominator come from
the same code path.

\subsection{Frozen result table}

The headline quantities of Section~V of the paper are collected in a
generated table that records, for each entry, its value, integer counts,
interval method, the exact denominator it is defined over, and the result
file it was read from. Conditional rates in this evaluation differ mainly in
their denominator---repair success over all unsafe first plans against unsafe
first plans on feasible intents, for instance---and the table exists so that
two such quantities cannot be interchanged. It is generated from the archived
results rather than transcribed, and it carries an entry recording how many
of its own source files were produced at a revision whose working tree was
not clean.
